% =====================================================================
%  Punim konference - formati IEEE (IEEEtran, conference)
%  Versioni i avancuar
%  Overleaf: ngarkoni main.tex + references.bib, kompiluesi: pdfLaTeX
% =====================================================================
\documentclass[conference]{IEEEtran}
\IEEEoverridecommandlockouts

\usepackage[T1]{fontenc}
\usepackage[utf8]{inputenc}
\usepackage[albanian]{babel}
\usepackage{cite}
\usepackage{amsmath,amssymb,amsfonts}
\usepackage{graphicx}
\usepackage{textcomp}
\usepackage{xcolor}
\usepackage{booktabs}
\usepackage{array}
\usepackage{url}

% Emërtimet në shqip (babel i rivendos në \begin{document}, prandaj \addto)
\addto\captionsalbanian{%
  \renewcommand{\abstractname}{Abstrakt}%
  \renewcommand{\tablename}{Tabela}%
}
\renewcommand{\IEEEkeywordsname}{Fjalë kyçe}

\begin{document}

% ---------------------------------------------------------------------
% 1. TITULLI I PUNIMIT
% ---------------------------------------------------------------------
\title{Modelet, korniza e mundësive dhe sfidat e eQeverisjes në Kosovë:\\
Një vlerësim i integruar i maturisë, kornizës institucionale\\
dhe koordinimit ndërinstitucional}

\author{
\IEEEauthorblockN{Emri Mbiemri}
\IEEEauthorblockA{\textit{Fakulteti / Departamenti} \\
\textit{Universiteti}\\
Prishtinë, Kosovë \\
email@domain.com}
}

\maketitle

% ---------------------------------------------------------------------
% ABSTRAKTI - do të plotësohet pasi të përfundojë punimi
% ---------------------------------------------------------------------
\begin{abstract}
[Do të plotësohet.]
\end{abstract}

% ---------------------------------------------------------------------
% 2. FJALËT KYÇE
% ---------------------------------------------------------------------
\begin{IEEEkeywords}
eQeverisja, ndërveprueshmëria, koordinimi ndërinstitucional
\end{IEEEkeywords}

% ---------------------------------------------------------------------
% 3. HYRJE
% ---------------------------------------------------------------------
\section{Hyrje}
\label{sec:hyrje}

Transformimi digjital i sektorit publik ka evoluar nga një iniciativë
kryesisht teknologjike në një proces të thellë reformimi institucional.
Nëse valët e para të eQeverisjes u përqendruan në publikimin e
informacionit dhe në digjitalizimin e formularëve, paradigma bashkëkohore
e \textit{qeverisjes së epokës digjitale} kërkon riintegrimin e
funksioneve të fragmentuara, organizimin e shërbimeve sipas nevojave të
qytetarit dhe ridizajnimin e proceseve administrative nga fillimi deri në
fund \cite{dunleavy2006,mergel2019}. Në kontekstin evropian, ky orientim
është institucionalizuar përmes parimeve si \textit{digjitale si
parazgjedhje}, \textit{vetëm një herë} (\textit{once-only}) dhe
ndërveprueshmëria si parakusht për shërbime publike të integruara
\cite{eif2017}.

Për Kosovën, kjo agjendë ka rëndësi të dyfishtë. Nga njëra anë, vendi ka
parakushte demografike dhe infrastrukturore relativisht të favorshme:
një popullsi prej rreth 1.6 milionë banorësh, ndër më të rejat në Evropë,
dhe një shkallë shumë të lartë të qasjes së ekonomive familjare në
internet \cite{ask2023}. % VERIFIKO: shifrat e fundit të ASK-së
Nga ana tjetër, procesi i integrimit evropian, i nisur formalisht me
hyrjen në fuqi të Marrëveshjes së Stabilizim-Asociimit në vitin 2016 dhe i
thelluar me aplikimin për anëtarësim në BE në dhjetor 2022, e shndërron
përafrimin me standardet evropiane të administratës publike digjitale në
një detyrim strategjik dhe jo vetëm në një zgjedhje politike
\cite{ec2024,sigma2023}.

Gjatë dekadës së fundit, Kosova ka ndërtuar elementet themelore të një
ekosistemi të eQeverisjes. Agjencia e Shoqërisë së Informacionit (ASHI)
është përcaktuar si autoriteti qendror për infrastrukturën shtetërore të
TIK-ut me Ligjin Nr.~04/L-145 \cite{ligji04L145}; platforma eKosova
funksionon si pikë e vetme qasjeje në shërbimet elektronike për qytetarët
dhe bizneset; platforma shtetërore e ndërveprueshmërisë synon shkëmbimin
e të dhënave ndërmjet regjistrave bazë; ndërsa Ligji Nr.~06/L-082 për
Mbrojtjen e të Dhënave Personale e ka përafruar kornizën kombëtare me
Rregulloren e Përgjithshme të BE-së për Mbrojtjen e të Dhënave (GDPR)
\cite{ligji06L082}.

Megjithatë, ekzistenca e këtyre elementeve nuk garanton integrimin e tyre
në një model koherent qeverisjeje. Literatura ndërkombëtare tregon se
shumica e projekteve të eQeverisjes në vendet në zhvillim dhe në tranzicion
dështojnë tërësisht ose pjesërisht, kryesisht për shkak të hendekut
ndërmjet projektimit formal dhe realitetit institucional \cite{heeks2003}.
Në rastin e Kosovës, ky hendek shfaqet potencialisht në tri dimensione:
(i) \textit{horizontalisht}, përmes zhvillimit të sistemeve sektoriale nga
ministri dhe agjenci të veçanta, pa standarde të përbashkëta semantike dhe
organizative; (ii) \textit{vertikalisht}, përmes asimetrisë ndërmjet
nivelit qendror dhe 38 komunave, të cilat sipas Ligjit Nr.~03/L-040 për
Vetëqeverisjen Lokale kanë kompetenca vetanake në ofrimin e shërbimeve
\cite{ligji03L040}, por kapacitete digjitale shumë të ndryshme; dhe
(iii) \textit{nga ana e kërkesës}, ku përdorimi i shërbimeve elektronike
varet nga besimi i qytetarëve, dobia e perceptuar dhe aftësitë digjitale
\cite{carter2005}.

Ky punim nis nga premisa se sfida kryesore e eQeverisjes në Kosovë nuk
është më mungesa e teknologjisë, por mungesa e koordinimit institucional
që e shndërron teknologjinë në vlerë publike. Kjo premisë mbështetet në
kornizën e \textit{zbatimit të teknologjisë} (\textit{technology
enactment}) të Fountain-it, sipas së cilës efektet e TIK-ut në qeveri
ndërmjetësohen nga format organizative dhe rregullimet institucionale
ekzistuese \cite{fountain2001}.

Pavarësisht rëndësisë së temës, ekziston një boshllëk i qartë njohurish.
Meqë Kosova nuk është shtet anëtar i Kombeve të Bashkuara, ajo nuk
përfshihet në Indeksin e Zhvillimit të eQeverisjes (EGDI)
\cite{undesa2024}, gjë që e privon vendin nga një krahasim i standardizuar
ndërkombëtar. Studimet ekzistuese për Kosovën janë kryesisht përshkruese,
sektoriale ose të lidhura me projekte të donatorëve, dhe rrallë e lidhin
vlerësimin e maturisë së shërbimeve me analizën e kornizës ligjore dhe me
dinamikat e qeverisjes shumënivelëshe. Në rrafsh teorik, modelet klasike të
maturisë janë kritikuar për supozimin e një progresi linear
\cite{coursey2008} dhe për mungesë ndjeshmërie ndaj kontekstit specifik të
vendeve në tranzicion \cite{janowski2015}. Për të adresuar këtë boshllëk,
punimi udhëhiqet nga pyetjet kërkimore (PK) në vijim:

\begin{IEEEdescription}[\IEEEsetlabelwidth{PK4:}]
  \item[PK1:] Në cilën fazë maturie ndodhet modeli aktual i eQeverisjes në
        Kosovë, në dimensionin vertikal (qendër--komunë) dhe horizontal
        (ndërsektorial)?
  \item[PK2:] Në çfarë mase korniza ligjore dhe institucionale ekzistuese
        mundëson apo kufizon integrimin dhe ndërveprueshmërinë e
        shërbimeve?
  \item[PK3:] Cilët faktorë të brendshëm dhe të jashtëm e dallojnë
        zhvillimin e eQeverisjes në nivel qendror nga ai në nivel lokal?
  \item[PK4:] Cilët mekanizma të qeverisjes së koordinimit mund ta
        reduktojnë fragmentimin institucional në proceset e digjitalizimit?
\end{IEEEdescription}

Kontributi i punimit është i trefishtë. Në rrafsh \textit{empirik}, ofrohet
një vlerësim i integruar i eQeverisjes në Kosovë, i cili ndërthur maturinë
e shërbimeve, kornizën normative dhe dimensionin qendër--lokal. Në rrafsh
\textit{teorik}, testohet zbatueshmëria e modeleve të maturisë dhe e qasjes
së kontekstualizimit në rastin e një shteti në tranzicion që mbetet jashtë
sistemit të matjeve të OKB-së. Në rrafsh \textit{praktik}, propozohet një
strategji e strukturuar e koordinimit ndërinstitucional, e mbështetur në
shtresat e Kornizës Evropiane të Ndërveprueshmërisë.

% ---------------------------------------------------------------------
% 4. QËLLIMI I PUNIMIT
% ---------------------------------------------------------------------
\section{Qëllimi i punimit}
\label{sec:qellimi}

Qëllimi i këtij punimi është të zhvillojë dhe të zbatojë një kuadër
vlerësimi të integruar për eQeverisjen në Kosovë, i cili lidh tri nivele
analize: (i) maturinë e shërbimeve elektronike, (ii) kapacitetin e
kornizës ligjore dhe institucionale për të mbështetur integrimin, dhe
(iii) mekanizmat e koordinimit ndërmjet institucioneve qendrore dhe
komunave. Në vend që të trajtohet si projekt teknologjik i izoluar,
eQeverisja konceptualizohet si sistem socio-teknik, rezultatet e të cilit
varen nga ndërveprimi i politikave, menaxhimit dhe teknologjisë
\cite{pardo2012,fountain2001}.

Kuadri analitik ndërthur modelet e maturisë të Layne dhe Lee
\cite{layne2001} dhe të Andersen dhe Henriksen \cite{andersen2006} me katër
shtresat e ndërveprueshmërisë (ligjore, organizative, semantike dhe
teknike) të Kornizës Evropiane të Ndërveprueshmërisë \cite{eif2017}. Në
rrafsh metodologjik, punimi mbështetet në një dizajn me metoda të
përziera, të strukturuar si studim rasti shumënivelësh, i cili përfshin:
analizë cilësore të dokumenteve (legjislacioni, strategjitë dhe raportet e
monitorimit të KE-së dhe SIGMA-s); vlerësim të strukturuar të një mostre
shërbimesh në platformën eKosova dhe në portalet e 38 komunave; si dhe
intervista gjysmë të strukturuara me zyrtarë të ASHI-t, të ministrive të
linjës dhe të komunave. Gjetjet sintetizohen përmes analizës SWOT dhe
zgjerohen në matricën TOWS, në mënyrë që diagnoza të shndërrohet në
opsione strategjike. Në fund, punimi synon t'u ofrojë politikëbërësve një
bazë empirike dhe një udhërrëfyes të zbatueshëm për kalimin nga
digjitalizimi i fragmentuar drejt shërbimeve publike të integruara dhe të
orientuara nga qytetari.

% ---------------------------------------------------------------------
% 5. OBJEKTIVAT E PUNIMIT
% ---------------------------------------------------------------------
\section{Objektivat e punimit}
\label{sec:objektivat}

\begin{enumerate}
  \item \textbf{O1 (PK1):} Të vlerësohet niveli i maturisë së shërbimeve
        elektronike në nivel qendror dhe lokal, përmes një indeksi të
        përshtatur nga modelet e maturisë.
  \item \textbf{O2 (PK2):} Të identifikohen boshllëqet, mbivendosjet dhe
        mospërputhjet e kornizës ligjore e institucionale në raport me
        \textit{acquis} të BE-së dhe me EIF-in.
  \item \textbf{O3 (PK3):} Të kryhet një analizë krahasuese SWOT për
        nivelin qendror dhe atë lokal dhe të zgjerohet në matricë TOWS.
  \item \textbf{O4 (PK4):} Të propozohet një model qeverisjeje për
        koordinimin ndërinstitucional, me role, instrumente dhe tregues
        të qartë të performancës.
\end{enumerate}

% ---------------------------------------------------------------------
% 6. GJENDJA AKTUALE (STATE OF THE ART)
% ---------------------------------------------------------------------
\section{Gjendja aktuale e kërkimeve të ngjashme}
\label{sec:sota}

\subsection{Modelet e maturisë dhe kritika e tyre}
Linja e parë e kërkimit e konceptualizon eQeverisjen si proces evolutiv me
faza. Modeli i Layne dhe Lee \cite{layne2001} dallon katër faza
(katalogimin, transaksionin, integrimin vertikal dhe integrimin
horizontal) dhe argumenton se me secilën fazë rritet kompleksiteti
teknologjik e organizativ, por edhe vlera për qytetarin. Andersen dhe
Henriksen \cite{andersen2006} e kritikojnë këtë qasje si tepër të
përqendruar te teknologjia dhe propozojnë modelin PPR, ku maturia matet
sipas shkallës së rindërtimit të proceseve dhe orientimit ndaj
përdoruesit. Janowski \cite{janowski2015} e zgjeron më tej këtë
perspektivë me një model evolutiv katërfazor, faza më e avancuar e të
cilit, kontekstualizimi, kërkon që politikat digjitale të përshtaten me
nevojat specifike të vendit, sektorit dhe komunitetit. Megjithatë, Coursey
dhe Norris \cite{coursey2008}, bazuar në të dhëna empirike nga qeveritë
lokale në SHBA, tregojnë se modelet e fazave nuk e parashikojnë mirë
zhvillimin real: shumica e institucioneve mbeten për një kohë të gjatë në
fazat informative dhe transaksionale. Rrjedhimisht, modelet e maturisë
duhet të përdoren si instrumente diagnostikuese dhe jo si trajektore
normative. Tabela~\ref{tab:modelet} paraqet një sintezë krahasuese të
modeleve kryesore dhe relevancën e tyre për rastin e Kosovës.

\begin{table*}[t]
\caption{Sintezë krahasuese e modeleve kryesore të maturisë së eQeverisjes
dhe relevanca e tyre për Kosovën}
\label{tab:modelet}
\centering
\footnotesize
\renewcommand{\arraystretch}{1.25}
\begin{tabular}{>{\raggedright\arraybackslash}p{2.5cm}
                >{\raggedright\arraybackslash}p{4.3cm}
                >{\raggedright\arraybackslash}p{3.6cm}
                >{\raggedright\arraybackslash}p{5.2cm}}
\toprule
\textbf{Modeli} & \textbf{Fazat} & \textbf{Fokusi analitik} &
\textbf{Relevanca për Kosovën} \\
\midrule
Layne dhe Lee \cite{layne2001} &
Katalogim; transaksion; integrim vertikal; integrim horizontal &
Integrimi teknologjik dhe funksional &
Bazë për matjen e integrimit qendër--komunë dhe ndërmjet sektorëve \\
Andersen dhe Henriksen \cite{andersen2006} &
Kultivim; zgjerim; maturim; revolucionarizim (PPR) &
Rindërtimi i proceseve dhe orientimi ndaj qytetarit &
Vlerëson nëse digjitalizimi shkon përtej portalit eKosova drejt
ridizajnimit të proceseve \\
Janowski \cite{janowski2015} &
Digjitalizim; transformim; angazhim; kontekstualizim &
Evolucioni institucional dhe kontekstual &
Përshtatja me kontekstin e tranzicionit, decentralizimit dhe integrimit
evropian \\
UN DESA -- EGDI \cite{undesa2024} &
Indeks i përbërë: shërbimet online (OSI), infrastruktura e
telekomunikacionit (TII), kapitali njerëzor (HCI) &
Krahasueshmëria ndërkombëtare &
Kosova është jashtë matjes; lind nevoja për vetëvlerësim me metodologji
ekuivalente \\
\bottomrule
\end{tabular}
\end{table*}

\subsection{Faktorët institucionalë të suksesit}
Linja e dytë fokusohet te përcaktuesit e suksesit. Gil-García dhe Pardo
\cite{gilgarcia2005} i grupojnë faktorët kritikë në pesë kategori:
informacioni dhe të dhënat, teknologjia, aspektet organizative e
menaxheriale, korniza ligjore e rregullatore, si dhe konteksti
institucional e mjedisor. Heeks \cite{heeks2003} vlerëson se rreth 35\% e
projekteve të eQeverisjes në vendet në zhvillim janë dështime totale dhe
rreth 50\% dështime të pjesshme, dhe këtë e shpjegon me hendekun ndërmjet
projektimit dhe realitetit në shtatë dimensione (modeli ITPOSMO). Fountain
\cite{fountain2001} shton se teknologjia nuk prodhon vetvetiu transformim:
ajo ``zbatohet'' përmes strukturave burokratike dhe rrjeteve ekzistuese, të
cilat shpesh e riprodhojnë fragmentimin në vend që ta eliminojnë.

\subsection{Ndërveprueshmëria dhe koordinimi ndërinstitucional}
Linja e tretë, më e lidhura drejtpërdrejt me këtë punim, trajton integrimin
dhe ndërveprueshmërinë. Scholl dhe Klischewski \cite{scholl2007} e
konceptualizojnë integrimin si qëllim organizativ dhe ndërveprueshmërinë si
aftësi teknike e semantike për ta realizuar atë, duke identifikuar
kufizime politike, ligjore, organizative dhe financiare. Pardo, Nam dhe
Burke \cite{pardo2012} tregojnë empirikisht se ndërveprueshmëria është
rezultat i ndërveprimit të politikave, menaxhimit dhe teknologjisë, ku
qeverisja e përbashkët dhe besimi ndërinstitucional janë po aq vendimtarë
sa standardet teknike. Korniza Evropiane e Ndërveprueshmërisë
\cite{eif2017} e operacionalizon këtë qasje përmes katër shtresave
(ligjore, organizative, semantike, teknike) dhe një shtrese mbikëqyrëse të
qeverisjes së shërbimeve publike të integruara. Rasti i Estonisë, shpesh i
cituar si pikë referimi, tregon se suksesi lidhet më shumë me kulturën e
inovacionit, vendimet e hershme arkitekturore dhe bashkëpunimin
publik--privat sesa me investimet teknologjike në vetvete
\cite{kalvet2012}.

\subsection{Dimensioni lokal dhe boshllëku kërkimor}
Në nivel lokal, Moon \cite{moon2002} gjen se adoptimi i eQeverisjes në
komuna mbetet kryesisht në faza të hershme, se retorika e transformimit e
tejkalon realitetin dhe se madhësia e kapacitetet administrative të
komunës ndikojnë ndjeshëm. Kjo perspektivë është veçanërisht relevante për
Kosovën, ku dallimet ndërmjet komunave të mëdha urbane dhe atyre të vogla
rurale janë të theksuara. Në sintezë, literatura ofron instrumente të
fuqishme për analizën e maturisë, të faktorëve të suksesit dhe të
ndërveprueshmërisë, por këto linja rrallë integrohen në një kuadër të
vetëm dhe zbatohen rrallë në kontekstin e shteteve në tranzicion me status
ndërkombëtar specifik, si Kosova. Ky punim synon ta mbushë pikërisht këtë
boshllëk.

% ---------------------------------------------------------------------
% SEKSIONET E TJERA - do të plotësohen në fazën e ardhshme
% ---------------------------------------------------------------------
\section{Vlerësimi i modelit aktual të eQeverisjes në Kosovë: gjendja dhe
rekomandimet}
\label{sec:vleresimi}
[Do të plotësohet.]

\section{Korniza ligjore dhe institucionale për zhvillimin e eQeverisjes
në Kosovë}
\label{sec:korniza}
[Do të plotësohet.]

\section{Analiza SWOT e zhvillimit të eQeverisjes në nivel qendror dhe
lokal}
\label{sec:swot}
[Do të plotësohet.]

\section{Strategji për përmirësimin e koordinimit ndërinstitucional në
proceset e digjitalizimit}
\label{sec:strategjia}
[Do të plotësohet.]

\section{Përfundime}
\label{sec:perfundime}
[Do të plotësohet.]

% ---------------------------------------------------------------------
% 7. REFERENCAT
% ---------------------------------------------------------------------
\bibliographystyle{IEEEtran}
\bibliography{references}

\end{document}


% =====================================================================
%  Referencat - stili IEEEtran
%  Referencat 1-14: punime shkencore të rishikuara nga kolegët
%  Referencat e tjera: dokumente institucionale, ligjore dhe statistikore
% =====================================================================

% ---------------- Modelet e maturisë ----------------
@article{layne2001,
  author  = {Layne, Karen and Lee, Jungwoo},
  title   = {Developing fully functional {E}-government: A four stage model},
  journal = {Government Information Quarterly},
  volume  = {18},
  number  = {2},
  pages   = {122--136},
  year    = {2001}
}

@article{andersen2006,
  author  = {Andersen, Kim Viborg and Henriksen, Helle Zinner},
  title   = {E-government maturity models: Extension of the {L}ayne and {L}ee model},
  journal = {Government Information Quarterly},
  volume  = {23},
  number  = {2},
  pages   = {236--248},
  year    = {2006}
}

@article{janowski2015,
  author  = {Janowski, Tomasz},
  title   = {Digital government evolution: From transformation to contextualization},
  journal = {Government Information Quarterly},
  volume  = {32},
  number  = {3},
  pages   = {221--236},
  year    = {2015}
}

@article{coursey2008,
  author  = {Coursey, David and Norris, Donald F.},
  title   = {Models of e-government: Are they correct? {A}n empirical assessment},
  journal = {Public Administration Review},
  volume  = {68},
  number  = {3},
  pages   = {523--536},
  year    = {2008}
}

% ---------------- Faktorët e suksesit dhe institucionet ----------------
@article{gilgarcia2005,
  author  = {Gil-Garc{\'i}a, J. Ram{\'o}n and Pardo, Theresa A.},
  title   = {E-government success factors: Mapping practical tools to theoretical foundations},
  journal = {Government Information Quarterly},
  volume  = {22},
  number  = {2},
  pages   = {187--216},
  year    = {2005}
}

@techreport{heeks2003,
  author      = {Heeks, Richard},
  title       = {Most e{G}overnment-for-Development Projects Fail: How Can Risks be Reduced?},
  institution = {Institute for Development Policy and Management, University of Manchester},
  type        = {iGovernment Working Paper},
  number      = {14},
  address     = {Manchester, UK},
  year        = {2003}
}

@book{fountain2001,
  author    = {Fountain, Jane E.},
  title     = {Building the Virtual State: Information Technology and Institutional Change},
  publisher = {Brookings Institution Press},
  address   = {Washington, DC},
  year      = {2001}
}

@article{dunleavy2006,
  author  = {Dunleavy, Patrick and Margetts, Helen and Bastow, Simon and Tinkler, Jane},
  title   = {New public management is dead---long live digital-era governance},
  journal = {Journal of Public Administration Research and Theory},
  volume  = {16},
  number  = {3},
  pages   = {467--494},
  year    = {2006}
}

@article{mergel2019,
  author  = {Mergel, Ines and Edelmann, Noella and Haug, Nathalie},
  title   = {Defining digital transformation: Results from expert interviews},
  journal = {Government Information Quarterly},
  volume  = {36},
  number  = {4},
  pages   = {101385},
  year    = {2019}
}

@article{carter2005,
  author  = {Carter, Lemuria and B{\'e}langer, France},
  title   = {The utilization of e-government services: Citizen trust, innovation and acceptance factors},
  journal = {Information Systems Journal},
  volume  = {15},
  number  = {1},
  pages   = {5--25},
  year    = {2005}
}

% ---------------- Integrimi dhe ndërveprueshmëria ----------------
@article{scholl2007,
  author  = {Scholl, Hans J. and Klischewski, Ralf},
  title   = {E-government integration and interoperability: Framing the research agenda},
  journal = {International Journal of Public Administration},
  volume  = {30},
  number  = {8--9},
  pages   = {889--920},
  year    = {2007}
}

@article{pardo2012,
  author  = {Pardo, Theresa A. and Nam, Taewoo and Burke, G. Brian},
  title   = {E-government interoperability: Interaction of policy, management, and technology dimensions},
  journal = {Social Science Computer Review},
  volume  = {30},
  number  = {1},
  pages   = {7--23},
  year    = {2012}
}

@article{kalvet2012,
  author  = {Kalvet, Tarmo},
  title   = {Innovation: A factor explaining e-government success in {E}stonia},
  journal = {Electronic Government, an International Journal},
  volume  = {9},
  number  = {2},
  pages   = {142--157},
  year    = {2012}
}

@article{moon2002,
  author  = {Moon, M. Jae},
  title   = {The evolution of e-government among municipalities: Rhetoric or reality?},
  journal = {Public Administration Review},
  volume  = {62},
  number  = {4},
  pages   = {424--433},
  year    = {2002}
}

% ---------------- Dokumente institucionale dhe ndërkombëtare ----------------
@misc{eif2017,
  author       = {{European Commission}},
  title        = {New {E}uropean Interoperability Framework: Promoting seamless services and data flows for {E}uropean public administrations},
  howpublished = {Publications Office of the European Union, Luxembourg},
  note         = {COM(2017) 134 final},
  year         = {2017}
}

@techreport{undesa2024,
  author      = {{United Nations Department of Economic and Social Affairs}},
  title       = {{United Nations E-Government Survey} 2024: Accelerating Digital Transformation for Sustainable Development},
  institution = {United Nations},
  address     = {New York, NY},
  year        = {2024}
}

@techreport{ec2024,
  author      = {{European Commission}},
  title       = {Kosovo 2024 Report},
  institution = {European Commission},
  type        = {Commission Staff Working Document},
  address     = {Brussels},
  year        = {2024}
}

@book{sigma2023,
  author    = {{OECD/SIGMA}},
  title     = {The Principles of Public Administration},
  publisher = {OECD Publishing},
  address   = {Paris},
  year      = {2023}
}

% VERIFIKO: titullin e saktë dhe vitin e botimit më të fundit të ASK-së
@techreport{ask2023,
  author      = {{Agjencia e Statistikave të Kosovës}},
  title       = {Përdorimi i teknologjisë së informacionit dhe komunikimit ({TIK}) në ekonomitë familjare},
  institution = {Agjencia e Statistikave të Kosovës},
  address     = {Prishtinë},
  year        = {2023}
}

% ---------------- Legjislacioni i Kosovës ----------------
@misc{ligji04L145,
  author       = {{Kuvendi i Republikës së Kosovës}},
  title        = {Ligji {Nr.}~04/{L}-145 për Organet Shtetërore të Shoqërisë së Informacionit},
  howpublished = {Gazeta Zyrtare e Republikës së Kosovës},
  year         = {2013}
}

@misc{ligji06L082,
  author       = {{Kuvendi i Republikës së Kosovës}},
  title        = {Ligji {Nr.}~06/{L}-082 për Mbrojtjen e të Dhënave Personale},
  howpublished = {Gazeta Zyrtare e Republikës së Kosovës},
  year         = {2019}
}

@misc{ligji03L040,
  author       = {{Kuvendi i Republikës së Kosovës}},
  title        = {Ligji {Nr.}~03/{L}-040 për Vetëqeverisjen Lokale},
  howpublished = {Gazeta Zyrtare e Republikës së Kosovës},
  year         = {2008}
}
